\section{泛化界的解释力与边界}

到目前为止得到的所有界都有一个共同签名：对最坏分布一致成立、常数保守、量级估计。一个诚实的问题必须现在提出：这些界在数值上几乎从来不具备预测力。本节把话说全——它们解释了什么（正则化、结构风险最小化、间隔），解释不了什么（深度网络、双下降），以及读者真正应该带走的判断框架。

\subsection{先承认：VC 界数值上很松}

代一组数字进去。\(\R^{100}\) 中的线性分类器有 \(d_{\mathrm{VC}}=101\)，取 \(n=10^5\)、\(\delta=0.05\)：

\[
\sqrt{\frac{d_{\mathrm{VC}}\log n+\log(1/\delta)}{n}}
\approx\sqrt{\frac{101\times11.5+3}{10^{5}}}
\approx 0.11.
\]

约 \(11\) 个百分点的界宽，还只是忽略常数后的量级；VC 定理原始形式中的常数会让它再宽几倍。对深度网络，VC 维至少随参数量增长（百万到十亿量级），代入任何现实的 \(n\)，界宽都超过 \(1\)——而错误率天生不超过 \(1\)，这样的界什么也没有说。

这不等于理论破产，而是定位不同：这些界对所有分布一致成立，为最坏情形付了全部保费。把``界大于 \(1\)''读作``理论错了''是误读；它的产出从来不在数值，而在结构——告诉我们泛化间隙随什么收紧、随什么变松。

\subsection{它解释了正则化与结构风险最小化}

把假设类组织成嵌套序列 \(\mathcal{H}_1\subset\mathcal{H}_2\subset\cdots\)，记 \(d_k=d_{\mathrm{VC}}(\mathcal{H}_k)\)，随 \(k\) 递增。对每个类各写一条界：

\[
R(h)\le\widehat R(h)+\mathrm{width}\bigl(d_k,n,\delta\bigr),
\qquad h\in\mathcal{H}_k.
\]

大类里 \(\widehat R\) 可以压得更低（偏差小），但界宽更大；小类相反。在 \(\mathcal{H}_k\) 上做 ERM 得到 \(\widehat h\)，把它的超出风险按类内最优 \(h_k^*\in\argmin_{h\in\mathcal{H}_k}R(h)\) 与全问题最优风险 \(R^*\) 拆开，两项各司其职：

\[
R\bigl(\widehat h\bigr)-R^*
=\underbrace{R\bigl(\widehat h\bigr)-R(h_k^*)}_{\text{估计误差}}
+\underbrace{R(h_k^*)-R^*}_{\text{逼近误差}}.
\]

类越大，第二项越小、第一项的界越宽。结构风险最小化（SRM）把这个权衡显式化：不再固定一个类，而是在嵌套序列上最小化``经验风险加界宽''，让数据自己挑选复杂度档位。惩罚型目标

\[
\min_{h}\;\widehat R(h)+\lambda\,\Omega(h)
\]

是它的连续版本：\(\Omega\) 大意味着``有效类''大，\(\lambda\) 就是为泛化间隙标出的价格。这解释了``光滑经验风险与 $L_2$ 正则化''中正则项的统计角色：它不只是改善条件数的数值装置，而是在为复杂度付费——偏差与界宽的权衡，被写进了目标函数。

\subsection{它解释了 SVM 为什么看间隔不看维度}

间隔界（margin bound）换了一个复杂度来源：若数据落在半径 \(B\) 的球内，且分类器以几何间隔 \(\gamma\) 分开数据，则有效复杂度的量级是 \((B/\gamma)^2\)，与输入维数无关——即使核方法把数据映到无穷维特征空间，界照常成立。``最大间隔与支持向量机''把目标写成``固定函数间隔、最小化 \(\norm{w}\)''，统计理由正在这里：在该约定下几何间隔是 \(1/\norm{w}\)，所以 SVM 的目标函数本身就是一条泛化界——hinge 项控制 \(\widehat R\)，\(\norm{w}^2\) 项控制界宽。

沿这条路走得更远的是范数界与 Rademacher 复杂度：用``类在这批数据上拟合随机噪声的平均能力''代替最坏情形的 VC 计数，常常给出紧得多的界，并且可以直接从数据估计。其推导与计算已构成一门课程，这里只登记它们的位置：当 VC 界太松而 margin 不适用时，它们是下一站。

\subsection{经典框架解释不了的部分}

也必须直面：现代深度学习的一些基本事实，经典一致收敛框架无法解释。

过参数化网络的参数量远超样本量，VC 维相应巨大；实验上，同样的网络甚至能以零训练误差拟合纯随机标签——也就是完全记住噪声。按经典理论，这样的类应当毫无泛化能力；然而同一架构、同一优化器在真实数据上泛化良好。问题出在哪一端，至今没有公认答案。

双下降（double descent）进一步撕开了经典图景：测试误差随模型规模先降后升，在插值阈值——模型容量刚好大到能把训练误差压到零的那一点——附近达到峰值。这一段还符合偏差与复杂度的权衡；但越过阈值继续增大模型，误差再次下降。一致的 U 形故事在插值点右侧整体失效。

共同的线索是：有效复杂度不只由假设类决定，还由优化过程决定。初始化、SGD 的隐式偏置、早停、数据本身的低维结构，都在把实际搜索范围收进 \(\mathcal{H}\) 的一个``好''子集里——这个子集的复杂度远小于整个类。隐式正则化、神经正切核、基于算法稳定性的界等方向都在改写``复杂度从哪来''的问题，但没有哪一个已经成为终局。这是开放问题，读者应把它当作前沿而非缺陷来记住。

\subsection{带走的判断框架}

泛化不是``模型简单''的同义词，而是``有效复杂度相对样本量受控''。控制的旋钮可以装在不同的地方：

\begin{itemize}
\tightlist
\item
  样本量：所有界都随 \(\sqrt{n}\) 收紧；
\item
  假设类大小：\(\log\abs{\mathcal{H}}\) 或 \(d_{\mathrm{VC}}\)；
\item
  几何量：间隔 \((B/\gamma)^2\)、范数 \(\norm{w}\)；
\item
  优化过程：隐式偏置——在经典框架之外，在现代实践之中。
\end{itemize}

这一整条理论链也把早已见过的工程纪律串成了必然推论：测试集只能用一次，调参必须经过验证集或重抽样（``交叉验证评估的是完整训练流程''），``训练 loss 下降''从来不是``会泛化''的证据。理论给出的不是免验证书，而是一张清单：每当数据参与一个决定，就问一次——这次决定花掉了多少随机性预算，它又对应界里的哪一项。
